% A single-column paper with an appendix of the kind coding-agent papers
% carry — made up, after "Automatic Generation of High-Performance RL
% Environments" (arXiv:2603.12145) — with what Reflow once read badly in
% one: headings in a bold sans set at the text's size, their numbers an em
% apart, straight under a table and its notes; numbered paragraphs led by a
% run-in heading, their lines turning over to the margin; listings of code
% between them; a prompt set in typewriter type, bold labels and all; and
% an algorithm ruled off as algorithmic's "ruled" style sets one.
% Rebuilt with: pdflatex agent-appendix.tex (TinyTeX: lm, booktabs, fancyvrb).
\documentclass[11pt]{article}
\usepackage[margin=1in]{geometry}
\usepackage[T1]{fontenc}
\usepackage{lmodern,amsmath,amssymb,booktabs,fancyvrb,float}
\makeatletter
\renewcommand\section{\@startsection{section}{1}{\z@}{-3ex \@plus -1ex}{1.5ex}{\normalfont\Large\sffamily\bfseries}}
\renewcommand\subsection{\@startsection{subsection}{2}{\z@}{-2.5ex \@plus -1ex}{1ex}{\normalfont\normalsize\sffamily\bfseries}}
\makeatother
\title{\textbf{Translating Environments with Coding Agents}}
\author{Ada Lindqvist \and Tomas Okafor}
\date{}
\begin{document}
\maketitle
\begin{abstract}
Environments for reinforcement learning are slow, and rewriting them by hand takes months. We have a coding agent translate them instead, and check each translation at four levels, from single functions to policies trained in one and evaluated in the other.
\end{abstract}
\section{Introduction}
Simulating the environment takes most of the time spent training an agent. Rewriting the environment for an accelerator takes a specialist months, and has been done for a handful of them. We let a coding agent do the rewriting, and check what it writes at four levels of increasing scope, so that the agent is told where a translation is wrong as well as that it is wrong. This paragraph runs on for a few lines so that the page has a measure of its body text, as a real paper's would.

The translations are as fast as the hand-written rewrites they are compared with, and pass the same tests. They cost a few dollars each, and the agent wrote every line of them.
\appendix
\section{Supplementary Details}
\subsection{Multi-Agent Validation}
We translated two of the environments again with two more agents, using the same prompts and the same tests. Both converge to translations that pass the tests (Table~\ref{tab:agents}), so the method does not depend on the agent.
\begin{table}[H]
\centering
\caption{\textbf{Multi-agent comparison.} Identical inputs; equivalent outputs.}\label{tab:agents}
\begin{tabular}{llrrr}
\toprule
\textbf{Environment} & \textbf{Agent} & \textbf{Iters} & \textbf{Tests} & \textbf{Cost} \\
\midrule
Pong & Agent One & 13 & 6/6 & \$0.05 \\
 & Agent Two & 3 & 5/6$^\S$ & \$0.08 \\
\midrule
HalfCheetah & Agent One & 20 & 69/69 & \$3.26 \\
 & Agent Three & 6 & 69/69 & --$^\dagger$ \\
\bottomrule
\end{tabular}

\vspace{4pt}
{\footnotesize $^\S$Same statistical test for both agents. $^\dagger$Cost not tracked.}
\end{table}
\subsection{Verification Ablation Details}
\textbf{HalfCheetah (complex physics).}\quad The run checked only by rollouts used eight end-to-end tests and forty-two iterations without converging. The agent could not tell a sign error in a force from a failure at the end of a rollout, and went round rewriting the vectorization instead.
\subsection{Cross-Hardware Validation}
\begin{table}[H]
\centering
\caption{\textbf{Throughput elsewhere.} Peak steps per second.}
\begin{tabular}{lrr}
\toprule
\textbf{Environment} & \textbf{SPS} & \textbf{vs. Reference} \\
\midrule
HalfCheetah & 13.1M & 290$\times$ \\
Pong & 1.4B & 23k$\times$ \\
\bottomrule
\end{tabular}
\end{table}
\subsection{Agent Translation Metrics}
The translation was carried out entirely through logged sessions. The first phase translated five core modules, and the second used a command-line agent for the six modules heaviest in logic.
\section{Performance Optimization Guide}
The following patterns, distilled from the case studies, make an environment faster. They are ordered by how much they usually help.

\medskip\noindent\textbf{1. Fixed-size state arrays.}\quad The compiler needs every shape known in advance. Replace all dynamic-length data structures (lists, dictionaries with varying keys) with fixed-size \texttt{jnp.ndarray} fields padded to their largest size, and mark the unused slots with a sentinel value.

\medskip\noindent\textbf{2. Branchless conditionals with \texttt{jnp.where}.}\quad Replace Python \texttt{if}/\texttt{else} with \texttt{jnp.where(condition, true\_val, false\_val)}, which computes both branches and selects the result by mask; this is faster on a GPU because it avoids divergence:
\begin{Verbatim}[xleftmargin=1em]
ball_vy = jnp.where(wall_hit, -ball_vy, ball_vy)
step = jax.vmap(
    partial(step, terrain=terrain),
    in_axes=(0, 0)
)
\end{Verbatim}
\medskip\noindent\textbf{3. JIT the outer interface.}\quad Apply \texttt{jax.jit} to the vectorized step and reset functions so that the whole batch is compiled to a single kernel, and compile them once while the environment is set up.
\subsection{Optimization Agent Prompt}
The following prompt is used once the environment passes the tests. It asks the agent to make it faster without changing what it does.
\begin{Verbatim}[xleftmargin=1em,fontsize=\small,commandchars=\\\{\}]
The environment has passed all verification tests. Now optimize it
for maximum steps per second.
\textbf{Target:}  Maximize SPS while keeping every test passing.
\textbf{Constraints:}
  - All tests must continue to pass after optimization
  - Do not change the external API (step, reset, observation
    shapes)
\end{Verbatim}
\begin{figure}[H]
\hrule height 0.6pt
\vspace{3pt}
\noindent\textbf{Algorithm 1 Hierarchical translation and verification.}
\vspace{3pt}
\hrule height 0.4pt
\vspace{3pt}
\noindent\begin{tabular}{@{}r@{\hspace{6pt}}l@{}}
1: & \textbf{Phase 1: Module translation} \\
2: & \textbf{for} $k = 1$ \textbf{to} $K$ \textbf{do} \\
3: & \hspace{1em}$m'_k \gets \textsc{Agent}(m_k)$ \\
4: & \hspace{1em}\textbf{if} $\textsc{RunTests}(m'_k) = \textsc{Pass}$ \textbf{then} \\
5: & \hspace{2em}\textbf{break} \\
6: & \hspace{1em}\textbf{end if} \\
7: & \textbf{end for} \\
8: & \textbf{return} $E_\text{perf}$ \\
\end{tabular}
\vspace{3pt}
\hrule height 0.6pt
\end{figure}
Algorithm 1 sets out the loop the agent runs, module by module, until every test passes.
\end{document}
